% Fragment -- inclus de main.tex prin \input
\clearpage
\setpartlabel{themeLighthouse}{\faChartBar}{Faza 6: Lighthouse}
\setcounter{section}{0}
\phantomsection
\addcontentsline{toc}{part}{\protect\textbf{Faza 6 \textendash{} Lighthouse}}
\handout{Faza 6 \textendash{} Lighthouse}{Audit automatizat. Indicatori de performanta. Core Web Vitals.}

Lighthouse este un instrument de audit automatizat dezvoltat de Google, integrat in Chrome DevTools si in serviciul web PageSpeed Insights. Capitolul prezinta cele patru categorii evaluate, modul de interpretare a raportului, indicatorii \emph{Core Web Vitals} folositi de Google in algoritmul de ranking, si remedierile pentru problemele frecvente.

\bigskip

\begin{outcomes}
Dupa capitolul asta vei \textcommabelow{s}ti:
\begin{itemize}[itemsep=1pt]
  \item Ruleaza un audit Lighthouse pe propria pagina prin PageSpeed Insights sau DevTools.
  \item Cite\textcommabelow{s}te raportul \textcommabelow{s}i identifica problemele cu impact maxim pe scor.
  \item Aplica fix-uri uzuale: \code{alt}, \code{lang}, \code{width}/\code{height} pe imagini, \code{defer} pe script.
  \item Optimizeaza imaginile cu format modern (WebP), \code{srcset}, \code{loading="lazy"}.
  \item Interpreteaza Core Web Vitals (LCP, INP, CLS) \textcommabelow{s}i pragurile lor.
\end{itemize}
\end{outcomes}

\begin{whymatters}
Lighthouse este standardul industrial pentru evaluarea unei pagini web. Scorul 95+ pe toate cele 4 categorii este criteriul implicit la code review in firmele moderne. Mai mult, Core Web Vitals influenteaza direct ranking-ul Google. Un audit de 2 minute identifica 80\% din problemele care ar costa ore de cautare manuala.
\end{whymatters}

\begin{nota}[Pre-rechizite]
Capitolul presupune un site deployat (Faza 4) si configurat cu meta tags (Faza 5). Auditul ruleaza pe URL-ul public.
\end{nota}

% =========================================================
\section{Categoriile de evaluare}

Lighthouse genereaza un scor 0\textendash{}100 pe fiecare dintre cele patru axe:

\begin{itemize}
  \item \textbf{Performance} \textendash{} timpul de incarcare si receptivitatea paginii.
  \item \textbf{Accessibility} \textendash{} conformitatea cu standardele WCAG (in principal contrast, structura semantica, asocieri label\textendash{}input, atribute \code{alt}).
  \item \textbf{Best Practices} \textendash{} folosirea de protocoale si API-uri moderne, absenta erorilor in consola, lipsa codului depasit.
  \item \textbf{SEO} \textendash{} prezenta meta tag-urilor obligatorii, structura semantica, indexabilitatea.
\end{itemize}

Scorurile se interpreteaza astfel: peste 90 = bun; peste 95 = nivel profesional; sub 50 = necesita interventie.

\subsection{Lab data vs field data}

Lighthouse foloseste \emph{lab data}: un scenariu sintetic in conditii standardizate (CPU lent, retea simulata 4G slow). Scorul reflecta capacitatea pe un dispozitiv mediu, nu situatia reala a utilizatorilor. Pentru date din productie, PageSpeed Insights afiseaza si \emph{field data} extrase din Chrome User Experience Report (CrUX), agregate pe ultimele 28 de zile.

% =========================================================
\section{Modalitati de executie}

\subsection{PageSpeed Insights (interfata web)}

Cea mai accesibila metoda. Pasi: accesarea \url{https://pagespeed.web.dev}, lipirea URL-ului public al paginii, executarea analizei. Avantaje: nu necesita instalare; include datele reale de la utilizatori (CrUX) atunci cand acestea exista; raport partajabil prin URL.

\subsection{Chrome DevTools (local)}

Procedura: deschiderea paginii in Chrome, \code{F12} $\to$ tab \emph{Lighthouse}, selectarea categoriilor, executarea \emph{Analyze page load}. Avantaje: functioneaza pe paginile servite de pe \code{localhost}, fara necesitatea de deploy; permite testare rapida intre modificari.

\subsection{CLI (pentru integrare CI)}

Lighthouse este disponibil ca pachet npm. Folosit in pipeline-uri CI/CD pentru a bloca deploy-urile care scad scorul sub un prag stabilit.

% =========================================================
\section{Structura raportului}

Fiecare categorie include:

\begin{itemize}
  \item Scorul numeric, evidentiat color (rosu / portocaliu / verde).
  \item Lista problemelor identificate, grupate pe tip. Selectarea unei probleme afiseaza descrierea, impactul, si pasii de remediere.
  \item Lista verificarilor trecute cu succes.
  \item Sectiunea \emph{Opportunities} (pentru Performance) listeaza imbunatatirile cu economia estimata in milisecunde.
\end{itemize}

% =========================================================
\section{Probleme frecvente si remediere}

\subsection{Performance}

\begin{itemize}
  \item \textbf{Imagini fara dimensiuni explicite.} Setarea atributelor \code{width} si \code{height} pe \code{<img>} permite browserului sa rezerve spatiu, prevenind \emph{Cumulative Layout Shift}.
  \item \textbf{Resurse care blocheaza randarea.} Adaugarea atributului \code{defer} pe \code{<script>} permite parsarea HTML in paralel.
  \item \textbf{Imagini supradimensionate.} O imagine 2400\,x\,1600 redata la 800 pixeli iroseste banda. Optimizare: \url{https://squoosh.app}.
  \item \textbf{Format imagini neoptimizat.} JPEG/PNG in locul WebP / AVIF.
  \item \textbf{Cache lipsa.} Resursele statice (CSS, JS, imagini) ar trebui servite cu header \code{Cache-Control: max-age=31536000} (1 an). Pe GitHub Pages: cache-ul este configurat automat, fara interventie.
\end{itemize}

\subsection{Accessibility}

\begin{itemize}
  \item \textbf{Atribut \code{alt} lipsa pe \code{<img>}.} Toate imaginile relevante semantic necesita \code{alt}; imaginile pur decorative folosesc \code{alt=""}.
  \item \textbf{Buton fara nume accesibil.} Includerea de text in interior sau adaugarea \code{aria-label}.
  \item \textbf{Etichete lipsa pe formulare.} Asocierea \code{<label for>} cu \code{<input id>}.
  \item \textbf{Contrast insuficient.} Conform WCAG AA, raportul de contrast minim este 4.5:1 pentru text normal si 3:1 pentru text mare. Verificare: \url{https://webaim.org/resources/contrastchecker/}.
\end{itemize}

\subsection{Best Practices}

\begin{itemize}
  \item \textbf{Erori in consola.} Lighthouse penalizeaza orice eroare \code{console.error} la incarcare. Verificare: \code{F12 $\to$ Console}.
  \item \textbf{Resurse incarcate prin HTTP (nu HTTPS).} Mixed content; rezolvat automat pe GitHub Pages.
\end{itemize}

\subsection{SEO}

\begin{itemize}
  \item \textbf{Document fara \code{<title>}.}
  \item \textbf{Lipsa \code{<meta name="description">}.}
  \item \textbf{Element \code{<html>} fara \code{lang}.}
  \item \textbf{Linkuri fara text descriptiv} (\code{<a href="X">click aici</a>} este penalizat).
\end{itemize}

% =========================================================
\section{Core Web Vitals}

Cei trei indicatori folositi explicit de Google in algoritmul de ranking, incepand cu 2021:

\begin{itemize}
  \item \textbf{LCP} (\emph{Largest Contentful Paint}) \textendash{} timpul pana la afisarea celui mai mare element vizibil. Prag: $<$2.5 secunde.
  \item \textbf{INP} (\emph{Interaction to Next Paint}) \textendash{} latenta intre o actiune a utilizatorului si raspunsul vizual. Prag: $<$200 milisecunde.
  \item \textbf{CLS} (\emph{Cumulative Layout Shift}) \textendash{} masura instabilitatii vizuale in timpul incarcarii. Prag: $<$0.1.
\end{itemize}

% =========================================================
\section{Optimizarea imaginilor}

Imaginile reprezinta, in medie, 50\% din dimensiunea unei pagini web. Optimizarea lor are impact direct asupra scorului Performance.

\subsection{Format modern}

WebP reduce dimensiunea cu 25\textendash{}35\% fata de JPEG, AVIF cu 50\%. Conversia se face cu \url{https://squoosh.app} (interfata web, fara instalare).

\subsection{Atribute necesare pe \texttt{<img>}}

\begin{dano}[Asa nu]
\begin{lstlisting}[style=html, numbers=none]
<img src="hero.jpg">
\end{lstlisting}
\textbf{Probleme.} Lipseste \code{alt} (a11y); lipsesc \code{width}/\code{height} (CLS); lipseste \code{loading} (download eager); o singura sursa pentru toate ecranele (banda risipita pe mobil).
\end{dano}

\begin{dado}[Asa da \textendash{} optimizat complet]
\begin{lstlisting}[style=html, numbers=none]
<img src="hero-800.webp"
     srcset="hero-400.webp 400w,
             hero-800.webp 800w,
             hero-1600.webp 1600w"
     sizes="(min-width: 768px) 50vw, 100vw"
     width="800" height="450"
     loading="lazy"
     alt="Echipa BEST Iasi in fata sediului UAIC">
\end{lstlisting}
\textbf{Beneficii.} Browserul alege rezolutia potrivita din \code{srcset}; imaginile sub fold sunt incarcate doar la nevoie; rezerva spatiu prin \code{width}/\code{height}; alt asigura accesibilitatea.
\end{dado}

\subsection{Lazy loading}

Atributul \code{loading="lazy"} amana descarcarea imaginilor aflate sub \emph{fold} (zona vizibila initial). Suportul browser este universal incepand cu 2020.

\textbf{Exceptie:} pentru imaginea \emph{hero} (prima imagine din \emph{viewport}), se foloseste \code{loading="eager"} sau atributul lipseste; \code{loading="lazy"} aici intarzie LCP-ul.

% =========================================================
\section{Imbunatatiri suplimentare pentru Performance}

\subsection{Preload pentru resurse critice}

Resursele critice (font principal, imaginea hero) pot fi declarate cu \code{<link rel="preload">} pentru ca browserul sa le descarce inainte de a parsa CSS-ul.

\begin{lstlisting}[style=html]
<link rel="preload" href="fonts/lato.woff2" as="font" type="font/woff2" crossorigin>
<link rel="preload" href="hero.webp" as="image">
\end{lstlisting}

\subsection{CSS critic inline}

Pentru paginile cu prima vizita prioritara (landing pages), CSS-ul critic (necesar pentru randarea zonei vizibile) poate fi inclus direct in \code{<head>} prin \code{<style>}, restul fiind incarcat asincron. Tehnica nu este necesara pentru pagini personale; merita atentie pentru site-uri comerciale.

% =========================================================
\section*{Verificare cuno\textcommabelow{s}tin\textcommabelow{t}e}

\begin{selfcheck}
\begin{enumerate}[itemsep=2pt]
  \item Care sunt cele 4 categorii evaluate de Lighthouse?
  \item Diferen\textcommabelow{t}a intre \emph{lab data} \textcommabelow{s}i \emph{field data}?
  \item Care sunt cele 3 metrici Core Web Vitals \textcommabelow{s}i pragurile lor?
  \item \emph{Spot the issue:} \code{<img src="hero.jpg">} pe o pagina de productie.
  \item Cand este corect sa folose\textcommabelow{s}ti \code{loading="lazy"} \textcommabelow{s}i cand nu?
\end{enumerate}
\end{selfcheck}

% =========================================================
\section*{Recapitulare}

\begin{recap}[Faza 6 -- Lighthouse]
\begin{itemize}[itemsep=1pt]
  \item Patru categorii: Performance, Accessibility, Best Practices, SEO.
  \item Praguri: 90+ acceptabil, 95+ profesional, sub 50 necesita interventie.
  \item Lab data (Lighthouse) vs field data (CrUX); PageSpeed Insights le combina.
  \item Executie: \url{https://pagespeed.web.dev} (web) sau Chrome DevTools $\to$ \emph{Lighthouse}.
  \item Top fix-uri: \code{alt} pe imagini, \code{lang} pe \code{<html>}, contrast $\geq$4.5:1, \code{width}/\code{height} pe \code{<img>}, \code{defer} pe \code{<script>}.
  \item Core Web Vitals: LCP $<$2.5s, INP $<$200ms, CLS $<$0.1.
  \item Imagini: WebP/AVIF, \code{srcset}, \code{loading="lazy"} (cu exceptia imaginii hero), dimensiuni explicite.
\end{itemize}
\end{recap}

\partdivider

% =========================================================
\begin{joc}
\textbf{Platforma:} PageSpeed Insights (\url{https://pagespeed.web.dev}).

\medskip

\textbf{Activitatea.}
\begin{enumerate}[itemsep=1pt]
  \item \emph{Pasul 1 (3 minute).} Lipirea URL-ului paginii proprii deployate. Executarea analizei. Scorul tipic, dupa parcurgerea fazelor 1\textendash{}5, este peste 90.
  \item \emph{Pasul 2 (5 minute).} Identificarea uneia sau a doua probleme cu impact ridicat. Aplicarea fix-ului in cod local. \code{git push}. Reexecutarea analizei dupa 60 de secunde.
  \item \emph{Pasul 3 (2 minute).} Aplicarea aceluiasi audit pe varianta deployata a folder-ului \code{kit-starter/01-broken/}. Compararea scorurilor.
\end{enumerate}

\textbf{Continuare independenta:} folder-ul \code{01-broken} contine 10 probleme documentate in fisierul README. Tinta este atingerea scorului 95+ pe toate cele patru categorii.

\medskip

\textbf{Durata estimata:} 10 minute.
\end{joc}

% =========================================================
\section*{\faGraduationCap\ \ Resurse pentru aprofundare}

\begin{itemize}[itemsep=2pt]
  \item Google. \emph{PageSpeed Insights}. \url{https://pagespeed.web.dev}
  \item Google. \emph{Learn Performance}. \url{https://web.dev/learn/performance}
  \item Google. \emph{Web Vitals}. \url{https://web.dev/articles/vitals}
  \item Google. \emph{Lighthouse documentation}. \url{https://developer.chrome.com/docs/lighthouse/overview}
  \item WebAIM. \emph{Contrast Checker}. \url{https://webaim.org/resources/contrastchecker/}
  \item Google. \emph{Squoosh image optimizer}. \url{https://squoosh.app}
  \item Real Favicon Generator. \url{https://realfavicongenerator.net}
\end{itemize}

\bigskip

\bridge{Capitolul Resurse \& urmatorii pa\textcommabelow{s}i (cap. 7) propune un roadmap de 4 saptamani pentru consolidare \textcommabelow{s}i o lista curatata de tutoriale, comunita\textcommabelow{t}i \textcommabelow{s}i tool-uri pentru a continua singur.}
